% !TeX root = ../main.tex
\section{Results}
\label{sec:results}

The results are organized to avoid conflating two different questions. We first ask whether a low-precision model trained with analog perturbations transfers to a fresh hardware instance. We then reintroduce PCM device physics to determine which precision regime remains deployable after that transfer problem is separated from retention and readout limits.

\subsection{IdealDevice Ablation: Fixed-Mask Failure and Algorithmic Rescue}
\label{subsec:failure-rescue}

We first use the tested AIHWKit IdealDevice quantization/noise baseline as an algorithmic ablation, not as the physical PCM deployment result. The fixed-mask HAT protocol mirrors the common practice of optimizing against one sampled hardware instance or analog tile configuration \citep{joshi2020accurate,rasch2021aihwkit}. This setting is useful because it exposes the boundary of instance-specific training. As summarized in Table~\ref{tab:algorithmic_rescue}, standard per-batch hardware-aware training reaches 87.28\% fresh-instance accuracy at 8-bit precision, but falls to 14.64\% at 4-bit (Figure~\ref{fig:algorithmic_rescue}A). The failure is therefore not ordinary optimization failure; it is a transfer failure caused by fresh hardware instance shifts at extreme low precision.

\begin{table}[htbp]
\centering
\caption{\textbf{IdealDevice pure-quantization ablation.} Data show fresh-instance accuracy across three seeds for the AIHWKit IdealDevice model. These results establish algorithmic cross-instance robustness, not PCM physical deployment performance.}
\label{tab:algorithmic_rescue}
\begin{tabular}{lccl}
\toprule
\textbf{Model Configuration} & \textbf{Precision} & \textbf{Fresh Accuracy} & \textbf{Role} \\
\midrule
IdealDevice Baseline & 8-bit & 87.28 $\pm$ 0.13\% & Stable upper bound \\
IdealDevice Baseline & 4-bit & 14.64 $\pm$ 0.11\% & Generalization collapse \\
Ensemble HAT & 4-bit & 86.16 $\pm$ 0.19\% & Algorithmic rescue \\
\bottomrule
\end{tabular}
\end{table}

To localize the error source, we remove constraints sequentially (Table~\ref{tab:ablation_mapping}). Quantization and hybrid mapping alone introduce only a small penalty ($98.06\% \rightarrow 97.39\%$). The collapse appears when a fixed mismatch mask is introduced, indicating that the model has learned a solution tied to one hardware draw rather than a solution that transfers across the mismatch distribution. Ensemble HAT mitigates this failure by forcing the optimizer to see a sequence of mismatch fields during training.

\begin{table}[htbp]
\centering
\caption{\textbf{Ablation of mapping and noise constraints.} Accuracy reported for the Tiny-ViT-5M backbone at 4-bit precision (where applicable) on CIFAR-10.}
\label{tab:ablation_mapping}
\begin{tabular}{lcclc}
\toprule
\textbf{Configuration} & \textbf{Precision} & \textbf{Partition} & \textbf{Mismatch Mask} & \textbf{Accuracy} \\
\midrule
Digital Reference & 32-bit & Digital & None & 98.06 $\pm$ 0.14\% \\
Hybrid Baseline & 4-bit & Analog/Digital & None & 97.39 $\pm$ 0.08\% \\
Fixed-mask HAT & 4-bit & Analog/Digital & Single (fixed) & 10.00 $\pm$ 0.00\% \\
Ensemble HAT & 4-bit & Analog/Digital & Resampled & 86.16 $\pm$ 0.19\% \\
\bottomrule
\end{tabular}
\end{table}

By resampling the device-to-device (D2D) variation pattern during training (Figure~\ref{fig:algorithmic_rescue}A), Ensemble HAT restores this 4-bit pure-quantization ablation to 86.16\% fresh-instance accuracy. The interpretation is deliberately narrow: 4-bit operation can be algorithmically viable when device physics are idealized, but this does not imply that 4-bit PCM remains stable after programming. Loss-landscape interpolation along the D2D mismatch direction (Figure~\ref{fig:algorithmic_rescue}B) supports this interpretation. Ensemble HAT finds a broad plateau in the mismatch direction, whereas Standard HAT converges to a sharp, instance-specific minimum.

\subsection{Realistic PCM Training and the Precision-Retention Frontier}
\label{subsec:pcm-precision-ladder}

Having isolated the algorithmic transfer problem, we next ask whether a realistic memory model preserves the same low-precision promise. We evaluate the trained deployment with the \texttt{PCMPresetUnitCell} model in AIHWKit, which includes nonlinear update physics, finite dynamic range, and time-dependent conductance drift.

Under the tested PCM UnitCell simulation regime, training remains viable at 4-bit, 6-bit, and 8-bit precision, but viability is not monotonic in bit width. We mark a setting as drift-constrained when its 24-hour accuracy loss exceeds a 1~pp deployment SLA. By this criterion, 8-bit and 6-bit PCM are effectively drift-flat, whereas 4-bit PCM incurs a 4.01~pp retention penalty and is therefore drift-limited (Table~\ref{tab:pcm_precision_ladder}). The important exception to a simple ``more bits is better'' reading is the 6-bit row: it is drift-stable, yet it has the lowest mean fresh accuracy. This behavior is consistent with a quantization-D2D interaction in which mismatch perturbations become comparable to the 1/64 conductance step. At 4-bit, coarser quantization masks part of the D2D perturbation; at 8-bit, the finer grid absorbs it; at 6-bit, the two effects interact unfavorably in this tested regime.

\begin{table}[htbp]
\centering
\caption{\textbf{PCM Precision Ladder and Stability Frontier.} Fresh accuracy reflects corrected new-protocol evaluations (3 seeds for 4/8-bit; 4 seeds for 6-bit). The retention columns are from the drift-evaluation protocol; $\Delta$ drift is computed as retention-eval 0~s accuracy minus 24~h accuracy, not as fresh-eval mean minus 24~h accuracy. The 6-bit row uses the completed four-seed drift closure and shows drift stability but low mean fresh accuracy, consistent with a D2D-sensitive transition regime rather than a deployment frontier optimum.}
\label{tab:pcm_precision_ladder}
\begin{tabular}{lcccll}
\toprule
\textbf{Precision} & \textbf{Fresh Acc.} & \textbf{1~h Acc.} & \textbf{24~h Acc.} & \textbf{$\Delta_{0\rightarrow24h}$} & \textbf{Deployment Role} \\
\midrule
\textbf{8-bit PCM} & 77.60\% & 77.49\% & 77.57\% & 0.04~pp & Drift-flat reference \\
\textbf{6-bit PCM} & 68.44\% & 68.46\% & 68.36\% & 0.04~pp & D2D-sensitive transition zone \\
\textbf{4-bit PCM} & 76.68\% & 74.04\% & 72.64\% & 4.01~pp & Drift-limited regime \\
\bottomrule
\end{tabular}
\end{table}

\subsection{Iso-Accuracy Operating Envelope}
\label{subsec:iso-accuracy}

The same separation appears at the readout level. In the joint ADC-precision and D2D-variability sweep (Figure~\ref{fig:physical_frontier}), accuracy is near chance below 5-bit ADC, and moving from 5 to 6 bits produces a consistent $\sim$7~pp gain. Once readout precision reaches this threshold, the limiting factor changes: further degradation is governed mainly by D2D variability rather than ADC quantization. The Sobol analysis captures this hierarchy, with $S_{\mathrm{ADC}}=0.98$ over the full grid and $S_{\mathrm{D2D}}=0.92$ after restricting to the operational regime. The resulting hardware prescription is sequential rather than simultaneous: secure 6-bit readout first, then spend device-engineering effort on mismatch reduction.

\begin{figure}[htbp]
\centering
\makebox[\textwidth][c]{\includegraphics[width=1.0\textwidth]{fig2_physical_frontier}}
\caption{\textbf{Physical constraints and the iso-accuracy operating envelope.} Contour map under joint ADC precision and D2D mismatch sweep. Readout precision forms a hard cliff below 6 bits, above which D2D variability dictates the operational boundary. Bold contour lines indicate the 80\%, 85\%, and 88\% accuracy levels.}
\label{fig:physical_frontier}
\end{figure}

The operating envelope turns the ablations into a deployment sequence rather than a single best training recipe. First, cross-instance overfitting must be removed by training against a distribution of mismatch fields. Second, the readout interface must clear the 6-bit ADC cliff. Third, the memory precision must be selected under the target retention constraint. Under the tested PCM model, 8-bit PCM is the strongest operating point because it combines 77.60\% fresh accuracy with negligible drift. 4-bit PCM remains competitive immediately after programming (76.68\%) but loses $\sim$4.01~pp over one day and violates the 1~pp drift SLA. 6-bit PCM is not drift-limited, but its 68.44\% fresh accuracy and high seed-to-seed variance (std 6.03\%) mark it as a D2D-sensitive transition zone rather than a deployment optimum.
